\section{Experimental design}
\label{sec:experimental-design}

\subsection{Predictions}
\label{sec:design-hypotheses}

Nine predictions are tested. They follow from the specification of the test and from the propositions of Section~\ref{sec:theory}, each with the criterion that decides it; Table~\ref{tab:hypotheses} lists them with the measured values and the outcomes. H1 to H5 were fixed before any experiment was run. H6 to H9, and the additive-lengthening arms on which H6, H7 and H9 are tested, were formulated after a first run of the re-estimated experiment, restricted to the override of the weights, had been inspected, and the clause of H4 that re-runs the outlying cell under further seeds was added after the stated-parameter experiment; the predictions are derived from the propositions, but they were not made blind, and the table marks them. \textbf{H1}, the specification of the product: on the stored history of case C14 the test is not evaluated at 2021-Q4, which is a single date, red at 2022-Q2 with one flagged change and red at 2022-Q4 with two, and the core share and the core maturity that the case holds fixed are not flagged. \textbf{H2}, the specification again: the test is green for a parameter held through a breach, for a change with no breach and for a change before the breach, and red for a change at the breach, and the stored series with its breach flags set to false is green; Proposition~\ref{prop:exact} adds that the engine equals the closed-form count on every history. \textbf{H3}, from Proposition~\ref{prop:separation}: on a stated parameter the test flags every history in which an assumption was changed at a breach date, so that the lower end of the 95\,\% interval of its detection is at least 99\,\% in every cell. \textbf{H4}, from Proposition~\ref{prop:false-flag}: the false-flag rate on a stated parameter with unrelated changes equals $\E[1-(1-\lambda)^m]$. \textbf{H5}, the conventional expectation of a statistical test: that false-flag rate is at most 5\,\%; Proposition~\ref{prop:false-flag}(c) predicts that it fails, and the prediction is included because a validator will read the test as one. \textbf{H6}: on a re-estimated parameter the test separates a 2-year override at breach dates from none, with an area under the ROC curve of at least 0.70; Proposition~\ref{prop:degeneracy} predicts that it does not. \textbf{H7}, from Proposition~\ref{prop:conformal}: the tolerance rule calibrated at 5\,\% holds its size on held-out banks and flags at least 90\,\% of 4-year overrides. \textbf{H8}: the tolerance calibrated at a breach prevalence of 0.25 keeps the point estimate of its false-flag rate at most 5\,\% at the prevalences 0.10 and 0.50, which Proposition~\ref{prop:conformal} does not guarantee since the prevalence changes the law of the statistic; the criterion is the point estimate, and the interval and the spread of a single calibration are reported beside it. \textbf{H9}: the look-ahead variant detects a 4-year override made at the date before a breach begins (a date at which $b_t=0$ and $b_{t+1}=1$) at least 1.5 times as often as the same-date rule. The arm was first run with the override one date before a planted breach date; with a persistence of \designStay{} the chain is reversible, so that date is itself a breach date with probability \designStay{}, and the arm did not test an anticipated override. It was replaced by the arm described here before the results of this work were written, and the size at which H9 is decided, 4 years, is the one at which the tolerance has power. A prediction that fails is reported as failed, with its numbers.

\subsection{Data and its provenance}
\label{sec:design-data}

Three sources of histories are used. The first is the stored output of the case C14 of the product, the wholesale deposits of a bank whose call-report aggregates are real and whose behavioural assumptions are synthetic: the stored documents of the three valuation dates (31 December 2021, 30 June 2022 and 31 December 2022) and of the fitted models, which hold the fitted weights, the weighted average maturity, the requested core share and core maturity, and the outlier ratio and breach flag at each date. Their inputs are the quarterly call reports of the bank from the public FDIC BankFind service, the file for the certificate 24735 from 31 December 2012 to 31 December 2022, and the daily par yield curves of the US Treasury for 2012 to 2022; both are in the licence class that allows mirroring, and the files are identified in the stored documents by their SHA-256 digests (the call-report file begins \shaFdic{} and the 2022 curve file \shaTreasury{}). The raw files are held outside the repository; the experiments read them from the data root of the device, rebuild the rolling fits of the replicating portfolio with the pipeline's own function, and check that the fits at the three valuation dates equal the stored ones to within \cOneFourRebuiltMaxDiff{} years. The rebuilt series has \realFits{} quarterly fits from \realFirstFit{} to the end of 2022.

The second source is a simulated population of stated parameters and the third a simulated population of banks with re-estimated parameters; both are described here. A history of the stated population has $T=\designT{}$ quarters. A breach flag series is a two-state Markov chain with stationary prevalence $p$ and persistence $\Prob(b_t=1\mid b_{t-1}=1)=\designStay{}$, so that the probability of entering a breach is $p(1-\designStay{})/(1-p)$, redrawn until at least one breach follows the first date. The weighted average maturity of a simulated bank is the engine's replicating-portfolio fit on a circular block bootstrap of the real quarterly rows \citep{politis1992}, a resampling of blocks of consecutive rows \citep{kunsch1989}, with blocks of \blockLength{} quarters, eight burn-in rows that are dropped, and twelve consecutive fits on windows of 16 quarters. The bootstrap keeps the rows of the real bank, the deposit cost and the seven tranche yields of one quarter together, and the regime of 2012 to 2022 with it. A history of a stated parameter is built from a draw made once and shared by its two arms: it starts at 4 years, changes at each date with probability $\lambda$ by 0.25 to 1.0 years of random sign (floor at 0.5 years), and in the alternative arm is lengthened by 0.5 to 2.0 years at each planted date.

\subsection{Published methods compared}
\label{sec:design-comparators}

Every method is run on the same histories as the proposed test, and none enters through its reported numbers. The comparators are those that can read the same inputs. The breach-agnostic count flags a history when the parameter changed at least once, whatever the flags, and stands for the monitors of parameter stability that do not read the limit. The one-sided exact test of Fisher on the two-by-two table of change against breach \citep{fisher1935}, the one-sided Mann--Whitney test of the absolute changes at breach dates against those elsewhere \citep{mann1947}, and the one-sided Welch test of the signed changes at breach dates against those elsewhere \citep{welch1947} are run with the implementations of \texttt{scipy} at a size of 5\,\%; the Welch test is the unequal-variance form of the comparison that Chow's test makes at known dates \citep{chow1960}, and the breach-agnostic count stands for the monitors of parameter stability that do not read the limit. A test of a break at an unknown date \citep{andrews1993,bai1998} is run as the one-sided supremum-$t$ test of an upward shift in the level of the series, with at least three dates on each side of the break and the pooled variance of the two segments, in two forms: the test alone, which ignores the flags, and the break-then-coincide rule, which flags a history when the test finds a shift and the estimated first date of the new level is a breach date. Each is calibrated at 5\,\% by the conformal quantile of its statistic (zero for the second when the estimated date is not a breach date) on the same calibration histories as the tolerance rule. The cumulative sum \citep{page1954,brown1975} is not run as a detector of its own: it tests the same alternative, a shift in the mean, from the same partial sums, and it has no input for the flags. The engine's own test of the stability of the replicating-portfolio weights, which is red when the largest sup-norm distance between the weights of consecutive windows exceeds 0.25, is run on the fitted weights of the banks, where it applies, and on the stored weights of the case. The tolerance rule and its look-ahead variant are calibrated as in Algorithm~\ref{alg:calibration} on 1,000 banks that no test history shares, at 5\,\%; on a stated parameter the two level-shift rules are calibrated on 1,000 further null histories of each cell. The Welch and Mann--Whitney tests are undefined when the history has no date outside the breach or when all the changes are equal, and in that case the p-value is set to one so that the test does not flag.

\subsection{Metrics and statistics}
\label{sec:design-metrics}

The false-flag rate of a detector is the proportion of null histories it flags, and the detection rate is the proportion of alternative histories it flags, each as $\hat p=k/n$ with its 95\,\% Wilson interval \citep{wilson1927},
\begin{equation}
\frac{\hat p+z^2/(2n)\ \pm\ z\sqrt{\hat p(1-\hat p)/n+z^2/(4n^2)}}{1+z^2/n},\qquad z=1.96 .
\label{eq:wilson}
\end{equation}
Rates are also computed on each of five blocks of histories, and the standard deviation across the blocks, with $n-1$ in the denominator, is the measure of dispersion. The threshold-free comparison is the area under the ROC curve of the statistic a detector thresholds \citep{hanley1982},
\begin{equation}
\mathrm{AUC}=\frac{1}{n_1n_0}\sum_{i=1}^{n_1}\sum_{j=1}^{n_0}\Bigl(\indic[S_i^{(1)}>S_j^{(0)}]+\tfrac12\indic[S_i^{(1)}=S_j^{(0)}]\Bigr),
\label{eq:auc}
\end{equation}
with a percentile bootstrap interval from resamples of both arms \citep{efron1979}, 500 for the stated population and 300 for the re-estimated one. A p-value enters the area as its negative, and the statistic of the proposed test is its count, whose ties are counted as one half; with ties on both arms the area is below one for a detector that flags every alternative history but also flags most null ones. The agreement of the engine with the closed-form count is a count of disagreements, which has to be zero. A closed form is compared with a simulation by whether it lies inside the Wilson interval of the simulated rate, which for a 95\,\% interval misses about one comparison in twenty.

\subsection{Protocol, repetitions and seeds}
\label{sec:design-protocol}

The experiments are numbered by what they run, and each is a script of the repository that writes a JSON document under the manuscript folder.
\begin{enumerate}
\item \emph{Committed history and enumeration} (seed \seedRepro{} for the random histories): the test on the three stored dates and the three parameters of the wholesale deposits, the two other behavioural tests of the engine on the stored inputs, the eight constructed histories, every history of a binary grid of up to six dates (each pattern of changes and each pattern of flags), \randomN{} random histories with ties at a cap, and the effect of noise in the last digit of a held parameter.
\item \emph{Stated parameter} (seed \seedStated{}): \statedCells{} cells, nine of them the grid of $\lambda\in\{0.05,0.10,0.20\}$ and prevalence $p\in\{0.10,0.25,0.50\}$ at $\pi_m=0.5$, and two more that vary $\pi_m$ to 0.25 and to 1 at $\lambda=0.10$ and $p=0.25$; \statedN{} histories per arm in each cell, in \statedBlocks{} blocks of \statedPerBlock{}; the two arms share their breach flags and their unrelated changes.
\item \emph{Re-estimated parameter} (seed \seedFitted{}): \fitNCal{} calibration banks and \fitN{} test banks in \fitBlocks{} blocks of \fitPerBlock{}; breach prevalences 0.10, 0.25 and 0.50, with a calibration prevalence of \calPrev{}; an additive lengthening of 0.5, 1, 2 or 4 years and a movement of the weights toward the longest tranche of 5, 10, 20 or 40\,\% per planted date, each with probability $\pi_m=\piM{}$ at a breach date (at least one); an additive lengthening of 2 or 4 years from the date before the first date of a breach spell, at each such date with probability $\pi_m$ (at least once), on the banks that have such a date; the tolerance and the two level-shift rules calibrated at each prevalence and at 0.25 applied to all; and \realDraws{} draws of breaches and overrides on the real fitted series of the bank, its last twelve fits, with the tolerance calibrated on the simulated banks.
\item \emph{Checks of the propositions} (seed \seedTheory{}): \hyperReps{} placements for each of \hyperCells{} pairs $(k,m)$ against Equation~\eqref{eq:hypergeometric}; the closed form of Equation~\eqref{eq:false-flag} with the number of breach dates of the very histories of the stated cells, and under three further seeds for the cell that departs most; the equality of the Fisher and permutation p-values on \fisherPermN{} random histories by enumeration of the placements; Equation~\eqref{eq:gaussian-power} on \gaussReps{} Gaussian histories for $m=1,3,6$; and \confCalibrations{} calibrations of \confNcal{} maxima of three Gaussian changes each, evaluated on \confNtest{} fresh maxima.
\item \emph{Ablation}: the proposed test with each component removed or replaced, on the stated and the re-estimated population, by reading the detectors of the experiments above (Section~\ref{sec:ablation}).
\end{enumerate}
No experiment draws a random number except through the seeds of its script and the bootstrap of the area under the curve, whose generator is seeded from the same constant.
